El iode 131 ($^{131}$I), descobert per Glenn Seaborg i John Livingood el 1938, és un important radioisòtop que s'utilitza en la radioteràpia posterior a la tiroidectomia en els casos de càncer de tiroide. Té un període de semidesintegració de 8,02 dies i es transforma en xenó (Xe) mitjançant una emissió primària $\beta^-$, seguida d'una emissió $\gamma$ de 364~keV.

\begin{apartat}{1}
Escriviu les equacions nuclears corresponents als processos esmentats i calculeu el percentatge que quedarà d'una determinada quantitat inicial de $^{131}$I després de 24,06 dies.
\end{apartat}

\begin{apartat}{1}
Calculeu la longitud d'ona dels fotons $\gamma$.
\end{apartat}

\dades{%
  Constant de Planck, $h = 6,62\times 10^{-34}\ \mathrm{J\,s}$.\\
  Massa de l'electró $= 9,11\times 10^{-31}\ \mathrm{kg}$.\\
  Velocitat de la llum, $c = 3,00\times 10^{8}\ \mathrm{m\,s^{-1}}$.\\
  $1\ \mathrm{eV} = 1,60\times 10^{-19}\ \mathrm{J}$.\\
  Nombre atòmic del iode $= 53$.}
